% !TeX root = proceedings.tex
% Source: official theieie_conference_2column.docx, SHA256 47a1d41f8c90...
% Word points are PDF big points (bp). Word's 1.2-line rule was measured in
% native Word output: HY Shin Myeongjo 9 bp has a 14.04 bp baseline distance.
\usepackage{etoolbox}
\usepackage{xurl}
\usepackage[noadjust]{cite}
\let\originalcite\cite
\renewcommand{\citedash}{]-[}
\renewcommand{\citepunct}{,\ }
% The official DOCX's m:mathPr/m:mathFont is Cambria Math.
\usepackage{unicode-math}
% Other machines may lack this PC's local font configuration.
\providecommand{\proceedingstitlelatin}{\rmfamily}
\providecommand{\proceedingstitlehangul}{}
\providecommand{\proceedingsauthorlatin}{\rmfamily}
\providecommand{\proceedingsauthorhangul}{}
\urlstyle{same}
\DeclareMathSizes{9.03374}{9}{6}{5}
\DeclareMathSizes{9.03375}{9}{6}{5}
\IfFontExistsTF{Cambria Math}{%
  \setmathfont{Cambria Math}
  \setmathfont{Cambria Math}[range=it/{Latin,latin,Greek,greek}]
}{%
  \setmathfont{Latin Modern Math}
  \PackageWarningNoLine{ieie-fonts}{Cambria Math unavailable; using Latin Modern Math for editing}
}
% Use Unicode math alphabets throughout; legacy OML slots otherwise select
% upright ASCII glyphs from Cambria when the old Korean setup preserves them.
\AtBeginDocument{%
  \let\mathrm\symup\let\mathit\symit\let\mathcal\symcal\let\mathbb\symbb
  \begingroup
  \count0=65
  \loop\global\Umathcode\count0=7 \symoperators\numexpr"1D434+\count0-65\relax
    \advance\count0 by1 \ifnum\count0<91\repeat
  \count0=97
  \loop\global\Umathcode\count0=7 \symoperators\numexpr"1D44E+\count0-97\relax
    \advance\count0 by1 \ifnum\count0<123\repeat
  \global\Umathcode`h=7 \symoperators"210E
  \endgroup
  \protected\def\alpha{\Umathchar7 \symoperators"1D6FC}%
  \protected\def\gamma{\Umathchar7 \symoperators"1D6FE}%
  \protected\def\eta{\Umathchar7 \symoperators"1D702}%
  \protected\def\tau{\Umathchar7 \symoperators"1D70F}%
  \protected\def\omega{\Umathchar7 \symoperators"1D714}%
}
% Preserve the original script loss symbol and double-struck indicator in the
% Word math font, independently of legacy alphabet-slot remapping.
\newcommand{\paperLoss}{\Umathchar7 \symoperators"2112 }
\newcommand{\paperIndicator}{\Umathchar7 \symoperators"1D7D9 }
% Centered page numbers requested by the author; footer is outside body margins.
\setlength{\footskip}{28.35bp}
\makeatletter
\def\ps@plain{\let\@mkboth\@gobbletwo\let\@oddhead\@empty\let\@evenhead\@empty
  \def\@oddfoot{\hfil{\rmfamily\mdseries\fontsize{9bp}{14.04bp}\selectfont\thepage}\hfil}
  \let\@evenfoot\@oddfoot}
\makeatother
\pagestyle{plain}
\renewcommand{\paperbodyfont}{\rmfamily\mdseries\fontsize{9bp}{14.04bp}\selectfont}
\renewcommand{\papercaptionfont}{\rmfamily\mdseries\fontsize{9bp}{14.04bp}\selectfont}
\renewcommand{\paperTableFont}{\rmfamily\mdseries\fontsize{9bp}{14.04bp}\selectfont}
\renewcommand{\papersectionfont}{\rmfamily\mdseries\fontsize{12bp}{18.72bp}\selectfont}
\renewcommand{\papersubsectionfont}{\rmfamily\mdseries\fontsize{10bp}{15.6bp}\selectfont}
\renewcommand{\paperreferencefont}{\rmfamily\mdseries\fontsize{9bp}{14.04bp}\selectfont}
\setlength{\parindent}{6bp} % Two literal spaces in the sample body paragraphs.
\setlength{\parskip}{0bp}
\setlength{\columnsep}{28.35bp}
\setlength{\emergencystretch}{1em}
% The local Korean body font maps bold to the regular face; synthesize weight
% explicitly for headings so Korean headings are visibly bold too.
\titleformat{\section}{\centering\papersectionfont\bfseries\addhangulfontfeatures{FakeBold=2}}{\Roman{section}.}{.5em}{}
\titleformat{\subsection}{\papersubsectionfont\bfseries\addhangulfontfeatures{FakeBold=2}}{\thesubsection}{.5em}{}
% The sample separates section headings from prose with one empty body line.
\titlespacing{\section}{0bp}{14.04bp}{14.04bp}
\titlespacing{\subsection}{0bp}{14.04bp}{0bp}
\captionsetup{
  font=paperCaption,labelfont=normalfont,textfont=normalfont,
  labelsep=period,justification=centering,singlelinecheck=false,
  skip=3bp
}
% Table captions sit directly above the top rule, without an empty body line.
\captionsetup[table]{skip=0bp}
\renewcommand{\arraystretch}{1}
\setlength{\textfloatsep}{14.04bp}
\setlength{\floatsep}{14.04bp}
\setlength{\intextsep}{14.04bp}
% One body-line gap outside each display block; no extra gap between its rows.
% Reapply locally because size switches can reset LaTeX's display skips.
\newcommand{\paperdisplayspacing}{%
  \setlength{\abovedisplayskip}{14.04bp}%
  \setlength{\belowdisplayskip}{14.04bp}%
  \setlength{\abovedisplayshortskip}{14.04bp}%
  \setlength{\belowdisplayshortskip}{14.04bp}%
  \setlength{\jot}{0bp}%
}
\AtBeginDocument{\paperdisplayspacing}
\AtBeginEnvironment{equation}{\paperdisplayspacing}
\AtBeginEnvironment{align}{\paperdisplayspacing}
% Fix the visible space outside the complete mathematical block, rather than
% adding font-dependent interline glue around individual display rows.
\makeatletter
\newbox\paperdisplaybox
% Optional optical offset balances visible ink above/below without changing
% the display's total height or the spacing between its mathematical rows.
\NewDocumentEnvironment{paperdisplay}{O{0bp}+b}{%
  % Allow a column break after the preceding complete paragraph.
  % Otherwise its final line can be dragged into the next column by the display.
  \par\vskip14.04bp\nointerlineskip
  \setbox\paperdisplaybox=\vbox{%
    \hsize=\columnwidth
    \setlength{\abovedisplayskip}{0bp}%
    \setlength{\belowdisplayskip}{0bp}%
    \setlength{\abovedisplayshortskip}{0bp}%
    \setlength{\belowdisplayshortskip}{0bp}%
    \setlength{\jot}{0bp}%
    \let\paperdisplayspacing\relax
    #2%
  }%
  \directlua{%
    local boxnum=\number\paperdisplaybox
    local box=tex.box[boxnum]
    local head=box.head
    local glue=node.id("glue")
    local penalty=node.id("penalty")
    local kern=node.id("kern")
    local hlist=node.id("hlist")
    while head and (head.id==glue or head.id==penalty or head.id==kern or
      (head.id==hlist and head.height==0 and head.depth==0)) do
      local empty=head
      head=node.remove(head,empty)
      node.flush_node(empty)
    end
    box.head=nil
    local correction=tex.sp("#1")
    if correction>0 or correction<0 then
      local pad=node.new("kern")
      pad.kern=correction
      head=node.insert_before(head,head,pad)
      local compensate=node.new("kern")
      compensate.kern=-correction
      head=node.insert_after(head,node.tail(head),compensate)
    end
    tex.box[boxnum]=node.vpack(head)
  }%
  \box\paperdisplaybox
  \nobreak\vskip14.04bp\nointerlineskip
  \@endpetrue
}{}
\makeatother
% Keep every figure/table at a column top, in source order.
\floatplacement{figure}{!t}
\floatplacement{table}{!t}
\renewcommand{\topfraction}{.9}
\renewcommand{\textfraction}{.1}
\setcounter{topnumber}{3}
\setcounter{totalnumber}{3}
% Keep inline [n] attached to the preceding text; superscripts belong to the review.
\renewcommand{\cite}[1]{\unskip\nobreak\originalcite{#1}}
\renewcommand{\refname}{참고문헌}
\makeatletter
\patchcmd{\thebibliography}{\leftmargin\labelwidth}{\leftmargin=12.4bp}{}{
  \PackageError{ieie-format}{Bibliography indent patch failed}{Check the document class.}}
\patchcmd{\thebibliography}{\advance\leftmargin\labelsep}{\advance\leftmargin3bp\labelsep=3bp}{}{}
\apptocmd{\thebibliography}{
  \setlength{\itemsep}{0bp}\setlength{\parsep}{0bp}\setlength{\parskip}{0bp}
  \setlength{\labelsep}{3bp}
  \paperreferencefont
}{}{}
\makeatother
